\section{About the Committee}

\subsection{Mandate}

The Disarmament and International Security Committee (DISEC) is the First Committee of the United Nations General Assembly, and its authority flows directly from the UN Charter rather than from any standing resolution. Under Article 11(1) of Chapter IV, the General Assembly may consider the general principles of cooperation in maintaining international peace and security, including the principles governing disarmament and the regulation of armaments, and may make recommendations on those principles to members states, to the Security Council or both.\footnote{United Nations. (1945). \emph{Charter of the United Nations, Chapter IV: The General Assembly (Articles 9–22).}\url{https://www.un.org/en/about-us/un-charter/chapter-4}} DISEC is the body through which the Assembly exercises that function.

This produces a deliberately broad remit. The Committee considers every disarmament and international security question that fall within the scope of the charter or relates to the powers of other UN organs alongside the general principles of cooperation in maintaining peace and the regulation of armaments.\footnote{United Nations General Assembly. (n.d.). \emph{Disarmament and International Security (First Committee).} Retrieved June 14, 2026, from\url{https://www.un.org/en/ga/first/index.shtml}} In practice the agenda is organized into seven thematic clusters: nuclear weapons, other weapons of mass destruction, outer space (disarmament aspects), conventional weapons, other disarmament measures and international security, regional disarmament and security, and the disarmament machinery.\footnote{United Nations. (2024, October 3). \emph{First Committee, in brief organizational meeting, sets out guidelines for completing its work within allotted time frame} (GA/DIS/3734).\url{https://press.un.org/en/2024/gadis3734.doc.htm}} Each annual session runs in three stages: a general debate, thematic discussions across those clusters, and action on draft resolutions.\footnote{United Nations General Assembly. (n.d.). \emph{Disarmament and International Security (First Committee).} Retrieved June 14, 2026, from\url{https://www.un.org/en/ga/first/index.shtml}}

The defining limit on this mandate is that DISEC cannot compel anyone. Its resolutions are recommendations rather than a binding law. They carry weight because they are debated and adopted by all UN member states and can establish customs, standards, and normative expectations over time, but they hold no enforcement power of their own.\footnote{Reaching Critical Will. (n.d.). \emph{First Committee of the UN General Assembly.} Women's International League for Peace and Freedom. Retrieved June 14, 2026, from\url{https://www.reachingcriticalwill.org/disarmament-fora/unga}} The Charter draws the boundary more sharply still: under Article 12, while the Security Council is exercising its function over a dispute or situation, the GA may not make a recommendation on that matter unless the Council requests it.\footnote{United Nations. (1945). \emph{Charter of the United Nations, Chapter IV: The General Assembly (Articles 9–22).}\url{https://www.un.org/en/about-us/un-charter/chapter-4}} To conclude, DISEC sets norms and frameworks, it does not authorize force, impose sanctions, or adjudicate a live crisis the Council has seized.

DISEC also does not operate in isolation. It works in close coordination with the United Nations Disarmament Commission and the Geneva based Conference on Disarmament, and it receives substantive and organizational support from the United Nations Office for Disarmament Affairs (UNODA).\footnote{United Nations General Assembly. (n.d.). \emph{Disarmament and International Security (First Committee).} Retrieved June 14, 2026, from\url{https://www.un.org/en/ga/first/index.shtml}} UNODA was first established in 1982 on the recommendation of the Assembly’s second special session on disarmament, restructured in 1998, and given its current name in 2007. It supports the norm-setting work of the First Committee, the Disarmament Commission, and the Conference on Disarmament and maintains instruments such as the UN Register of Conventional Arms.\footnote{United Nations Office for Disarmament Affairs. (n.d.). \emph{About us.} Retrieved June 14, 2026, from\url{https://disarmament.unoda.org/en/who-we-are/about-unoda}} One mark of the Committee’s standing is procedural: under Rule 58(a) of the Assembly's rules of procedure, the First Committee is the only Main Committee entitled to verbatim records of its meetings, the same word-for-word coverage given to the Assembly plenary and the Security Council.\footnote{United Nations General Assembly. (n.d.). \emph{Rules of procedure of the General Assembly, IX. Records (Rules 58–59).} Retrieved June 14, 2026, from \url{https://www.un.org/en/ga/about/ropga/recds.shtml}; United Nations Department for General Assembly and Conference Management. (n.d.). \emph{Verbatim reporting.} Retrieved June 14, 2026, from\url{https://www.un.org/dgacm/en/content/verbatim-reporting}}

For this year’s topic, the mandate maps cleanly onto the disarmament half of the agenda. UNODA’s remit explicitly covers practical, post-conflict disarmament measures such as disarming and demobilising former combatants, which places the disarmament and reintegration of fighters, including children, within the First Committee’s competence.\footnote{Nuclear Threat Initiative. (2025, September 18). \emph{United Nations Office of Disarmament Affairs (UNODA).}\url{https://www.nti.org/education-center/treaties-and-regimes/un-office-of-disarmament-affairs/}} The flow of small arms and light weapons that arm child soldiers, and the frameworks governing their collection and destruction during DDR, are core conventional arms business for DISEC. What sits outside the committee's reach is the humanitarian and child-protection machinery handled elsewhere in the UN system. While DISEC can build the disarmament architecture and the norms around reintegration, it cannot run prosecutions or direct field protection operations. Keeping resolutions on the disarmament and security side of that line is what keeps them in order.

\subsection{History}

DISEC’s history follows one clear trajectory. It began by trying to contain the deadliest weapons ever built, and over time broadened to confront the everyday weapons and practices that drive most actual conflict.

The committee is as old as the UN itself. When the organisation was founded in 1945, the General Assembly was tasked with considering the principles of disarmament and the regulation of armaments, and the First Committee became the standing forum for that work.\footnote{United Nations General Assembly. (n.d.). \emph{Disarmament and International Security (First Committee).} Retrieved June 14, 2026, from\url{https://www.un.org/en/ga/first/index.shtml}; United Nations. (1945). \emph{Charter of the United Nations, Chapter IV: The General Assembly (Articles 9–22).}\url{https://www.un.org/en/about-us/un-charter/chapter-4}}, For its first decades, that meant the Cold War and the nuclear arms race. The debate of this era produced the foundation of the nonproliferation regime, the Treaty on the Non-Proliferation of Nuclear Weapons, opened for signature in 1968 and in force since 1970 (cite).

The system that still governs UN disarmament was set in 1978 at the General Assembly’s first special session devoted entirely to disarmament, whose Final Document defined the priorities that continue to guide this work.\footnote{United Nations Office for Disarmament Affairs. (n.d.). \emph{Treaty on the Non-Proliferation of Nuclear Weapons (NPT).} Retrieved June 14, 2026, from\url{https://disarmament.unoda.org/en/our-work/weapons-mass-destruction/nuclear-weapons/treaty-non-proliferation-nuclear-weapons}} That session built the wider machinery DISEC operates within, including the Disarmament Commission and the Geneva-based Conference on Disarmament.

The end of the Cold War changed the agenda. With the superpower standoff over, the committee expanded beyond weapons of mass destruction to the threats that shape conflict on the ground: the illicit trade in small arms and light weapons, the disarmament and demobilization of former combatants and, more recently, the military use of emerging technologies.\footnote{United Nations Office for Disarmament Affairs. (n.d.). \emph{About us.} Retrieved June 14, 2026, from\url{https://disarmament.unoda.org/en/who-we-are/about-unoda}} The clearest result of this shift is the Arms Trade Treaty adopted by the General Assembly in 2013 and in force since 2014, the first binding global rules on the trade in conventional weapons from small arms to warships.\footnote{Nuclear Threat Initiative. (2025, September 18). \emph{United Nations Office of Disarmament Affairs (UNODA).}\url{https://www.nti.org/education-center/treaties-and-regimes/un-office-of-disarmament-affairs/}}\footnote{United Nations Regional Centre for Peace and Disarmament in Asia and the Pacific. (n.d.). \emph{Arms Trade Treaty.} Retrieved June 14, 2026, from\url{https://www.unrcpd.org/conventional-weapons/arms-trade-treaty}} Norm building continued in 2017 with the Treaty on the Prohibition of Nuclear Weapons in force since 2021.\footnote{United Nations Office for Disarmament Affairs. (n.d.). \emph{Treaty on the Prohibition of Nuclear Weapons.} Retrieved June 14, 2026, from\url{https://disarmament.unoda.org/en/our-work/weapons-mass-destruction/nuclear-weapons/treaty-prohibition-nuclear-weapons}}
